\cotefile{162-5}
\scanpage{26}
Utilisons
\[
\boxed{\Theta_X\underset{\unsure{\triangleright}}{=}i_X\delta+\delta i_X}\quad\inserted{\text{(en degré quelconque)}}\quad\text{cf plus bas}
\]
il vient alors, au choix
\[
\begin{aligned}
T&=\sum i_X i_Y\delta i_Z-i_Z\delta i_X i_Y\deleted{\unsure{-\delta i_Z i_X i_Y}}\qquad \unsure{!|/}\quad\text{ou}\\
T(\delta\omega)&=\sum[i_{X\wedge Y},\Theta_Z](\delta\omega)
\end{aligned}
\]
\margin{pourquoi cela est-il nul ??}

On a donc une \uuline{$\Theta$-Algèbre de Lie} sur $X_0$, i.e. un \uline{Anneau de Lie} \deleted{sur} $\mathscr T$ sur $\deleted{\unsure{\underline O_{X_0}}}X_0$, muni d'un \uline{hom} \inserted{$\Theta:X\longmapsto\Theta_X$.} d'\uline{Anneaux} de Lie dans $\mathscr T_{X_0/\unsure{Z|\mathbb Z}}$, et d'une structure de $\underline O_{X_0}$-Module rendant $\Theta$ $\underline O_{X_0}$-linéaire et satisfaisant la formule $[\deleted{\unsure{\Theta}}X,fY]=\deleted{\unsure{f\Theta}}f[X,Y]+\Theta_X(f)Y$.

4.\quad \uline{Complexe de Cochaines associé à une $\Theta$-algèbre de Lie} \inserted{$\mathscr T$}

Cf. \unsure{mes notes|une note} Bourbaki ; on pose \deleted{(\unsure{supposant} qu'on soit sur un pt)}
\[
\underline{\unsure{C|\mathcal C}}^p=\underline{\mathrm{Hom}}(\bigwedge^p\mathscr T,O_{X_0})
\]
et on y définit une différentielle par une formule connue ..... .

Lorsque $\mathscr T$ provient d'un complexe de De Rham $(\bigwedge^*\Omega,\delta)$, alors l'hom d'alg. graduées évident $\bigwedge^*\Omega\longrightarrow C^*(\mathscr T)$ est \uline{compatible} avec les \deleted{structures} différentielles.

\margin{25}
